\chapter{Canonical Residual Closure and Syntactic Mealy Categories}
\label{ch:canonical-residual-closure-behaviour-specifications}

\section{Overview}
\label{sec:canonical-residual-closure-overview}

The preceding chapter constructed the final behaviour system
\[
\mathsf{Beh}_{\mathbb A,\mathbb B}
\]
for Mealy morphisms
\[
\mathbb A\rightsquigarrow\mathbb B,
\]
and showed that every Mealy system has a unique behaviour map into it. It
also identified minimal separated realizations with derivative closures and,
under exactness, with Nerode quotients.

This chapter develops the corresponding residual, recognizer-theoretic, and
syntactic theory. The basic object is not an arbitrary Mealy system, but a
behaviour specification: a subobject
\[
S\hookrightarrow \mathsf{Beh}_{\mathbb A,\mathbb B}
\]
in the slice over
\[
A_0\times B_0.
\]
A raw presentation
\[
k:K\to \mathsf{Beh}_{\mathbb A,\mathbb B}
\]
contributes only its image subobject. Thus the theory depends on the
behaviours specified, not on the chosen presentation of them.

From such a subobject we construct its canonical residual closure
\[
\mathsf D(S)\hookrightarrow \mathsf{Beh}_{\mathbb A,\mathbb B}.
\]
It is the image of the derivative-evaluation map
\[
A_1\ltimes S\to \mathsf{Beh}_{\mathbb A,\mathbb B},
\qquad
(a,\beta)\longmapsto \partial_a\beta.
\]
Since \(A_1\) is the full arrow object of the input internal category
\(\mathbb A\), this image already contains derivatives by all composite input
arrows. Consequently
\[
S\longmapsto \mathsf D(S)
\]
is a closure operator on the subobject lattice of the final behaviour system.
Its fixed points are precisely the sub-Mealy systems of
\[
\mathsf{Beh}_{\mathbb A,\mathbb B}.
\]

The residual closure is the state object of the canonical generated separated
recognizer of \(S\). The corresponding syntactic Mealy category is the action
category
\[
\operatorname{Syn}_{\mathbb A,\mathbb B}(S)
:=
\int_{\mathbb A}\mathsf D(S),
\]
equipped with its canonical input and output labellings
\[
I_S:
\operatorname{Syn}_{\mathbb A,\mathbb B}(S)
\to
\mathbb A^{\mathrm{op}},
\qquad
\Omega_S:
\operatorname{Syn}_{\mathbb A,\mathbb B}(S)
\to
\mathbb B^{\mathrm{op}}.
\]
This labelled input-discrete internal category is generated by the specified
behaviours and is behaviourally separated. It is the canonical generated
separated recognizer of the specification.

The chapter also constructs a residual-expression Mealy system
\[
\mathsf E(S)
\]
whose states are formal residual expressions
\[
(a,\beta),
\qquad
\beta\in S.
\]
Its quotient by behavioural equality is exactly
\[
\mathsf D(S).
\]
Passing to action categories gives an effective quotient presentation of the
syntactic Mealy category:
\[
\int_{\mathbb A}\mathsf E(S)
\twoheadrightarrow
\operatorname{Syn}_{\mathbb A,\mathbb B}(S).
\]
Thus the syntactic Mealy category is obtained by taking formal residual
expressions and identifying precisely those with the same behaviour.

In the one-object Set case, where
\[
\mathbb A=M,
\qquad
\mathbb B=N
\]
are monoids, the residual closure is
\[
\mathsf D(S)
=
\{\partial_r\lambda\mid r\in M,\ \lambda\in S\}.
\]
The syntactic transition-output monoid is the image of
\[
M
\longrightarrow
\operatorname{MealyEnd}_N(\mathsf D(S)),
\]
which sends an input element to its induced residual-state transformation
together with its output function. If \(M=I^\ast\) is free, this recovers the
classical residual-state construction and the transition monoid of the
minimal derivative automaton, with the Mealy-output component retained.

\section{Ambient setting and notation}
\label{sec:canonical-residual-closure-ambient-setting}

Throughout this chapter, assume that
\[
\mathcal E
\]
is a Grothendieck topos. Thus \(\mathcal E\) has the finite-limit,
local-cartesian-closed, regular, exact, and locally presentable structure used
below. In particular, every slice
\[
\mathcal E/Z
\]
is again regular and exact.

Let
\[
\mathbb A
=
(A_0,A_1,s_A,t_A,e_A,m_A)
\qquad\text{and}\qquad
\mathbb B
=
(B_0,B_1,s_B,t_B,e_B,m_B)
\]
be internal categories in \(\mathcal E\). We keep the composition convention
used in the preceding chapters: if
\[
a:u\to v,
\qquad
b:v\to w,
\]
then
\[
b\circ a:u\to w
\]
means first \(a\), then \(b\). Thus
\[
m_A(a,b)=b\circ a.
\]

Write
\[
Z:=A_0\times B_0.
\]
Let
\[
\mathsf{Beh}
:=
\mathsf{Beh}_{\mathbb A,\mathbb B}
\]
denote the final behaviour system for Mealy morphisms
\[
\mathbb A\rightsquigarrow\mathbb B.
\]
Its carrier lies over \(Z\):
\[
(p_{\mathsf{Beh}},q_{\mathsf{Beh}}):
\mathsf{Beh}\to A_0\times B_0.
\]

Since \(\mathsf{Beh}\) is a Mealy system, it has transition and output maps
\[
\delta_{\mathsf{Beh}}:
A_1\times_{t_A,A_0,p_{\mathsf{Beh}}}\mathsf{Beh}
\to
\mathsf{Beh},
\]
and
\[
\omega_{\mathsf{Beh}}:
A_1\times_{t_A,A_0,p_{\mathsf{Beh}}}\mathsf{Beh}
\to
B_1.
\]
The typing equations are
\[
p_{\mathsf{Beh}}\delta_{\mathsf{Beh}}
=
s_A\pi_1,
\]
\[
s_B\omega_{\mathsf{Beh}}
=
q_{\mathsf{Beh}}\delta_{\mathsf{Beh}},
\qquad
t_B\omega_{\mathsf{Beh}}
=
q_{\mathsf{Beh}}\pi_2.
\]
Thus
\[
\omega_{\mathsf{Beh}}(a,\beta):
q_{\mathsf{Beh}}(\partial_a\beta)\to q_{\mathsf{Beh}}(\beta)
\]
is an arrow of \(\mathbb B\).

The output typing is recorded by the commutative square
\[
\begin{tikzcd}[column sep=large]
A_1\times_{t_A,A_0,p_{\mathsf{Beh}}}\mathsf{Beh}
\arrow[r, "\omega_{\mathsf{Beh}}"]
\arrow[d, "{(q_{\mathsf{Beh}}\delta_{\mathsf{Beh}},\,q_{\mathsf{Beh}}\pi_2)}"']
&
B_1
\arrow[d, "{(s_B,t_B)}"]
\\
B_0\times B_0
\arrow[r, equal]
&
B_0\times B_0 .
\end{tikzcd}
\]
The transition typing is the separate equation
\[
p_{\mathsf{Beh}}\delta_{\mathsf{Beh}}=s_A\pi_1.
\]

\begin{definition}[Derivative]
\label{def:canonical-behaviour-derivative}
Let
\[
a:u\to v
\]
be an arrow of \(\mathbb A\), and let
\[
\beta\in\mathsf{Beh}
\]
be a generalized behaviour lying over \(v\). The \textbf{derivative} of
\(\beta\) by \(a\) is
\[
\partial_a\beta
:=
\delta_{\mathsf{Beh}}(a,\beta).
\]
It lies over \(u\).
\end{definition}

\begin{proposition}[Derivative laws]
\label{prop:canonical-derivative-laws}
For every behaviour \(\beta\) over \(v\),
\[
\partial_{e_A(v)}\beta=\beta.
\]
For composable arrows
\[
a:u\to v,
\qquad
b:v\to w,
\]
and every behaviour \(\beta\) over \(w\), one has
\[
\partial_{b\circ a}\beta
=
\partial_a(\partial_b\beta).
\]
\end{proposition}

\begin{proof}
These are precisely the unit and multiplication laws for the transition map
of the Mealy system
\[
\mathsf{Beh}.
\]
Since the input action is a right action, a composite
\[
b\circ a:u\to w
\]
is processed first by \(b\), then by \(a\). Thus
\[
\delta_{\mathsf{Beh}}(b\circ a,\beta)
=
\delta_{\mathsf{Beh}}(a,\delta_{\mathsf{Beh}}(b,\beta)),
\]
which is the displayed derivative identity.
\end{proof}

\section{Behaviour specifications as subobjects}
\label{sec:behaviour-specifications-subobjects}

The final behaviour system is behaviourally separated. Consequently, a
behaviour specification is determined by the subobject of behaviours that it
presents.

\begin{definition}[Behaviour specification]
\label{def:behaviour-specification-subobject}
A \textbf{behaviour specification} is a subobject
\[
S\hookrightarrow \mathsf{Beh}
\]
in the slice
\[
\mathcal E/Z.
\]
\end{definition}

\begin{definition}[Image specification of a raw presentation]
\label{def:image-specification}
Let
\[
k:K\to\mathsf{Beh}
\]
be a morphism in
\[
\mathcal E/Z.
\]
Its \textbf{image specification} is the image subobject
\[
\overline K\hookrightarrow\mathsf{Beh}
\]
in the regular epi--mono factorization
\[
K\twoheadrightarrow \overline K\hookrightarrow \mathsf{Beh}
\]
computed in
\[
\mathcal E/Z.
\]
\end{definition}

\[
\begin{tikzcd}[column sep=large]
K
\arrow[rr, "k"]
\arrow[dr, two heads]
&&
\mathsf{Beh}
\\
&
\overline K
\arrow[ur, hook']
&
\end{tikzcd}
\]

In what follows, all constructions are applied to behaviour subobjects
\[
S\hookrightarrow\mathsf{Beh}.
\]
A raw presentation
\[
k:K\to\mathsf{Beh}
\]
is first replaced by its image specification
\[
\overline K\hookrightarrow\mathsf{Beh}.
\]

\section{Derivative evaluation}
\label{sec:derivative-evaluation}

Let
\[
S\hookrightarrow\mathsf{Beh}
\]
be a behaviour specification. Write
\[
p_S:S\to A_0
\qquad\text{and}\qquad
q_S:S\to B_0
\]
for the composites with
\[
p_{\mathsf{Beh}}
\qquad\text{and}\qquad
q_{\mathsf{Beh}}.
\]

\begin{definition}[Residual-expression object]
\label{def:residual-expression-object-over-specification}
The \textbf{residual-expression object} of \(S\) is
\[
A_1\ltimes S
:=
A_1\times_{t_A,A_0,p_S}S.
\]
A generalized element is a pair
\[
(a,\beta)
\]
where
\[
a:u\to v
\]
is an arrow of \(\mathbb A\), and
\[
\beta\in S
\]
lies over \(v\).
\end{definition}

\[
\begin{tikzcd}
A_1\ltimes S
\arrow[r]
\arrow[d]
&
S
\arrow[d, "p_S"]
\\
A_1
\arrow[r, "t_A"']
&
A_0
\arrow[ul, phantom, "\lrcorner", very near start]
\end{tikzcd}
\]

\begin{definition}[Derivative evaluation]
\label{def:derivative-evaluation}
The \textbf{derivative evaluation map} of \(S\) is
\[
d_S:
A_1\ltimes S\to \mathsf{Beh}
\]
defined by
\[
d_S(a,\beta)=\partial_a\beta.
\]
\end{definition}

\begin{lemma}[Slice structure on residual expressions]
\label{lem:residual-expression-slice-structure}
The object
\[
A_1\ltimes S
\]
is regarded as an object of
\[
\mathcal E/Z
\]
through the composite
\[
A_1\ltimes S
\xrightarrow{d_S}
\mathsf{Beh}
\xrightarrow{(p_{\mathsf{Beh}},q_{\mathsf{Beh}})}
Z.
\]
With this structure, \(d_S\) is a morphism in
\[
\mathcal E/Z.
\]
\end{lemma}

\begin{proof}
By definition, the structure map of \(A_1\ltimes S\) over \(Z\) is the
composite
\[
A_1\ltimes S
\xrightarrow{d_S}
\mathsf{Beh}
\to
Z.
\]
Therefore \(d_S\) is a morphism over \(Z\). In generalized-element notation,
the base point of
\[
(a,\beta)
\]
is
\[
\bigl(s_A(a),q_{\mathsf{Beh}}(\partial_a\beta)\bigr).
\]
\end{proof}

\[
\begin{tikzcd}[column sep=large]
A_1\ltimes S
\arrow[r, "d_S"]
\arrow[dr, "{(p,q)_{A_1\ltimes S}}"']
&
\mathsf{Beh}
\arrow[d, "{(p_{\mathsf{Beh}},q_{\mathsf{Beh}})}"]
\\
&
Z
\end{tikzcd}
\]

\begin{remark}[The projection to \(S\)]
\label{rem:projection-not-slice-map}
The projection
\[
A_1\ltimes S\to S
\]
is generally not a morphism in
\[
\mathcal E/Z
\]
when \(A_1\ltimes S\) is given the slice structure of
Lemma~\ref{lem:residual-expression-slice-structure}. It is only a morphism
over \(A_0\) before derivative evaluation.

Indeed, a pair
\[
(a,\beta)
\]
with
\[
a:u\to v
\]
and \(\beta\) over \(v\) is sent by derivative evaluation to a behaviour over
\(u\), and the \(B_0\)-component may change from
\[
q_{\mathsf{Beh}}(\beta)
\]
to
\[
q_{\mathsf{Beh}}(\partial_a\beta).
\]
Thus the relevant slice map is
\[
d_S:A_1\ltimes S\to\mathsf{Beh}.
\]
\end{remark}

\section{The residual closure operator}
\label{sec:residual-closure-operator}

\begin{definition}[Residual closure]
\label{def:residual-closure}
The \textbf{residual closure} of a behaviour specification
\[
S\hookrightarrow\mathsf{Beh}
\]
is the image in
\[
\mathcal E/Z
\]
of the derivative evaluation map
\[
d_S:A_1\ltimes S\to\mathsf{Beh}.
\]
It is denoted
\[
\mathsf D(S)\hookrightarrow\mathsf{Beh}.
\]
Thus there is a regular epi--mono factorization
\[
A_1\ltimes S
\twoheadrightarrow
\mathsf D(S)
\hookrightarrow
\mathsf{Beh}
\]
in
\[
\mathcal E/Z.
\]
\end{definition}

\[
\begin{tikzcd}[column sep=large]
A_1\ltimes S
\arrow[rr, "d_S"]
\arrow[dr, two heads, "e_S"']
&&
\mathsf{Beh}
\\
&
\mathsf D(S)
\arrow[ur, hook, "j_S"']
&
\end{tikzcd}
\]

Equivalently, \(\mathsf D(S)\) is the regular image of the family of all
derivatives
\[
\partial_a\beta,
\]
where
\[
\beta\in S
\]
and \(a\) is a type-compatible input arrow. Since \(A_1\) is the full arrow
object of \(\mathbb A\), derivatives by composite input arrows are included
at once.

\begin{proposition}[Extensivity]
\label{prop:residual-closure-extensive}
For every behaviour specification
\[
S\hookrightarrow\mathsf{Beh},
\]
there is an inclusion
\[
S\leq \mathsf D(S)
\]
in
\[
\operatorname{Sub}_Z(\mathsf{Beh}).
\]
\end{proposition}

\begin{proof}
Define
\[
\eta_S:S\to A_1\ltimes S
\]
by
\[
\eta_S(\beta)=\bigl(e_A(p_S\beta),\beta\bigr).
\]
With the slice structure of
Lemma~\ref{lem:residual-expression-slice-structure}, the map \(\eta_S\) is a
morphism in \(\mathcal E/Z\), because
\[
d_S\eta_S(\beta)
=
\partial_{e_A(p_S\beta)}\beta
=
\beta.
\]
Thus the inclusion
\[
S\hookrightarrow\mathsf{Beh}
\]
factors through the image of \(d_S\). Hence
\[
S\leq \mathsf D(S).
\]
\end{proof}

\[
\begin{tikzcd}[column sep=large]
S
\arrow[r, "\eta_S"]
\arrow[dr, hook']
&
A_1\ltimes S
\arrow[d, "d_S"]
\arrow[r, two heads]
&
\mathsf D(S)
\arrow[dl, hook]
\\
&
\mathsf{Beh}
&
\end{tikzcd}
\]

\begin{proposition}[Monotonicity]
\label{prop:residual-closure-monotone}
If
\[
S\leq T
\]
as subobjects of
\[
\mathsf{Beh},
\]
then
\[
\mathsf D(S)\leq \mathsf D(T).
\]
\end{proposition}

\begin{proof}
The inclusion
\[
S\hookrightarrow T
\]
induces a morphism
\[
A_1\ltimes S\to A_1\ltimes T
\]
over \(A_0\). This induced map is also a morphism in \(\mathcal E/Z\), because
\[
d_T\circ(A_1\ltimes S\to A_1\ltimes T)=d_S.
\]
The derivative evaluation maps therefore form a commutative square
\[
\begin{tikzcd}
A_1\ltimes S
\arrow[r]
\arrow[d, "d_S"']
&
A_1\ltimes T
\arrow[d, "d_T"]
\\
\mathsf{Beh}
\arrow[r, equal]
&
\mathsf{Beh}.
\end{tikzcd}
\]
Hence the image of \(d_S\) is contained in the image of \(d_T\), and
\[
\mathsf D(S)\leq \mathsf D(T).
\]
\end{proof}

\begin{proposition}[Idempotence]
\label{prop:residual-closure-idempotent}
For every behaviour specification
\[
S\hookrightarrow\mathsf{Beh},
\]
one has
\[
\mathsf D(\mathsf D(S))=\mathsf D(S).
\]
\end{proposition}

\begin{proof}
By extensivity, it suffices to prove
\[
\mathsf D(\mathsf D(S))\leq \mathsf D(S).
\]

Let
\[
A_1\ltimes S
\xrightarrow{e_S}
\mathsf D(S)
\xrightarrow{j_S}
\mathsf{Beh}
\]
be the image factorization of
\[
d_S.
\]
Here \(A_1\ltimes S\) is regarded as an object of \(\mathcal E/Z\) through
\(d_S\). Consequently its \(A_0\)-component is
\[
p_{A_1\ltimes S}
=
p_{\mathsf{Beh}}d_S
=
s_A\pi_1,
\]
not the projection
\[
p_S\pi_2=t_A\pi_1
\]
inherited from its construction as a pullback over \(A_0\).

Since regular epimorphisms are stable under pullback, the morphism
\[
A_1\times_{t_A,A_0,p_{A_1\ltimes S}}(A_1\ltimes S)
\longrightarrow
A_1\times_{t_A,A_0,p_{\mathsf D(S)}}\mathsf D(S)
\]
induced by \(e_S\) is a regular epimorphism.

\[
\begin{tikzcd}
A_1\times_{A_0}(A_1\ltimes S)
\arrow[r, two heads]
\arrow[d]
&
A_1\times_{A_0}\mathsf D(S)
\arrow[d]
\\
A_1\ltimes S
\arrow[r, two heads, "e_S"']
&
\mathsf D(S)
\arrow[ul, phantom, "\lrcorner", very near start]
\end{tikzcd}
\]

A generalized element of the upper-left object is a triple
\[
(c,a,\beta)
\]
with
\[
a:u\to v,
\qquad
c:w\to u,
\qquad
\beta\in S
\]
lying over \(v\). Pulling back the derivative map
\[
d_{\mathsf D(S)}:
A_1\ltimes \mathsf D(S)\to\mathsf{Beh}
\]
along the displayed regular epimorphism gives the assignment
\[
(c,a,\beta)
\longmapsto
\partial_c(\partial_a\beta).
\]
By Proposition~\ref{prop:canonical-derivative-laws},
\[
\partial_c(\partial_a\beta)
=
\partial_{a\circ c}\beta.
\]
Since
\[
m_A(c,a)=a\circ c,
\]
this assignment factors through
\[
d_S:A_1\ltimes S\to\mathsf{Beh}
\]
via
\[
(c,a,\beta)\longmapsto (m_A(c,a),\beta).
\]
Therefore, after pullback along a regular epimorphism, the map
\[
d_{\mathsf D(S)}
\]
lands in
\[
\mathsf D(S)\hookrightarrow\mathsf{Beh}.
\]

In a regular category, pullback along a regular epimorphism reflects
containment of subobjects. Hence
\[
d_{\mathsf D(S)}
\]
itself factors through
\[
\mathsf D(S)\hookrightarrow\mathsf{Beh}.
\]
Therefore
\[
\mathsf D(\mathsf D(S))\leq \mathsf D(S).
\]
Together with extensivity, this gives
\[
\mathsf D(\mathsf D(S))=\mathsf D(S).
\]
\end{proof}

\begin{theorem}[Residual closure is a closure operator]
\label{thm:residual-closure-operator}
The assignment
\[
S\longmapsto \mathsf D(S)
\]
defines a closure operator on the subobject poset
\[
\operatorname{Sub}_Z(\mathsf{Beh}).
\]
That is, it is extensive, monotone, and idempotent.
\end{theorem}

\begin{proof}
This is the combination of
Propositions~\ref{prop:residual-closure-extensive},
\ref{prop:residual-closure-monotone}, and
\ref{prop:residual-closure-idempotent}.
\end{proof}

\section{Closed behaviour systems}
\label{sec:closed-behaviour-systems}

\begin{definition}[Closed behaviour system]
\label{def:closed-behaviour-system}
A behaviour subobject
\[
D\hookrightarrow\mathsf{Beh}
\]
is called a \textbf{closed behaviour system} if
\[
\mathsf D(D)=D.
\]
\end{definition}

Thus a closed behaviour system is a behaviour subobject containing all its
derivatives.

\begin{theorem}[Closed subobjects and sub-Mealy systems]
\label{thm:closed-iff-submealy}
For a subobject
\[
D\hookrightarrow\mathsf{Beh}
\]
in
\[
\mathcal E/Z,
\]
the following are equivalent:
\begin{enumerate}
\item \(D\) is closed:
\[
\mathsf D(D)=D.
\]
\item The derivative map
\[
d_D:
A_1\times_{t_A,A_0,p_D}D\to\mathsf{Beh}
\]
factors through
\[
D\hookrightarrow\mathsf{Beh}.
\]
\item The inclusion
\[
D\hookrightarrow\mathsf{Beh}
\]
is the carrier map of a sub-Mealy system of
\[
\mathsf{Beh}.
\]
\end{enumerate}
\end{theorem}

\begin{proof}
The image of
\[
d_D:
A_1\ltimes D\to\mathsf{Beh}
\]
is, by definition,
\[
\mathsf D(D).
\]
Hence
\[
\mathsf D(D)\leq D
\]
if and only if
\[
d_D
\]
factors through
\[
D\hookrightarrow\mathsf{Beh}.
\]
By extensivity,
\[
D\leq\mathsf D(D).
\]
Therefore
\[
\mathsf D(D)=D
\]
if and only if \(d_D\) factors through \(D\). This proves the equivalence of
(1) and (2).

If the derivative map factors through \(D\), then the transition map of
\(\mathsf{Beh}\) restricts to a map
\[
A_1\times_{t_A,A_0,p_D}D\to D.
\]
The output map restricts by restricting
\[
\omega_{\mathsf{Beh}}
\]
to
\[
A_1\times_{t_A,A_0,p_D}D.
\]
The Mealy unit and multiplication equations hold because they hold in
\[
\mathsf{Beh}.
\]
Thus \(D\) is a sub-Mealy system of \(\mathsf{Beh}\).

Conversely, if \(D\) is a sub-Mealy system of \(\mathsf{Beh}\), then its
transition map is the restriction of the transition map of
\[
\mathsf{Beh}.
\]
Hence \(D\) is closed under derivatives, so
\[
\mathsf D(D)=D.
\]
\end{proof}

\begin{corollary}
\label{cor:closed-behaviours-submealy}
Closed behaviour systems are precisely the sub-Mealy systems of the final
behaviour system
\[
\mathsf{Beh}_{\mathbb A,\mathbb B}.
\]
\end{corollary}

\section{Reflection into closed behaviour systems}
\label{sec:reflection-closed-behaviour-systems}

Let
\[
\operatorname{ClosedBeh}_{\mathbb A,\mathbb B}
\]
denote the full subposet of
\[
\operatorname{Sub}_Z(\mathsf{Beh})
\]
consisting of closed behaviour systems.

\[
\begin{tikzcd}[column sep=large]
\operatorname{ClosedBeh}_{\mathbb A,\mathbb B}
\arrow[r, hook, "J"]
&
\operatorname{Sub}_Z(\mathsf{Beh})
\arrow[l, bend left=35, "\mathsf D"]
\end{tikzcd}
\]

\begin{theorem}[Reflection into closed behaviour systems]
\label{thm:reflection-closed-behaviour-systems}
The inclusion
\[
J:
\operatorname{ClosedBeh}_{\mathbb A,\mathbb B}
\hookrightarrow
\operatorname{Sub}_Z(\mathsf{Beh})
\]
has left adjoint
\[
\mathsf D:
\operatorname{Sub}_Z(\mathsf{Beh})
\to
\operatorname{ClosedBeh}_{\mathbb A,\mathbb B}.
\]
Equivalently, for every behaviour specification
\[
S\hookrightarrow\mathsf{Beh}
\]
and every closed behaviour system
\[
D\hookrightarrow\mathsf{Beh},
\]
one has
\[
\mathsf D(S)\leq D
\quad\Longleftrightarrow\quad
S\leq D.
\]
\end{theorem}

\begin{proof}
Suppose first that
\[
\mathsf D(S)\leq D.
\]
Since
\[
S\leq \mathsf D(S)
\]
by extensivity, it follows that
\[
S\leq D.
\]

Conversely, suppose that
\[
S\leq D
\]
and \(D\) is closed. By monotonicity,
\[
\mathsf D(S)\leq\mathsf D(D).
\]
Since \(D\) is closed,
\[
\mathsf D(D)=D.
\]
Thus
\[
\mathsf D(S)\leq D.
\]
This proves the adjunction.
\end{proof}

\begin{definition}[Canonical residual realization]
\label{def:canonical-residual-realization}
For a behaviour specification
\[
S\hookrightarrow\mathsf{Beh},
\]
the closed behaviour system
\[
\mathsf D(S)\hookrightarrow\mathsf{Beh}
\]
is called the \textbf{canonical residual realization} of \(S\).
\end{definition}

\begin{proposition}[Minimality]
\label{prop:canonical-residual-realization-minimal}
The canonical residual realization
\[
\mathsf D(S)
\]
is the smallest sub-Mealy system of
\[
\mathsf{Beh}
\]
containing
\[
S.
\]
\end{proposition}

\begin{proof}
By Corollary~\ref{cor:closed-behaviours-submealy}, sub-Mealy systems of
\[
\mathsf{Beh}
\]
are precisely closed behaviour systems. By
Theorem~\ref{thm:reflection-closed-behaviour-systems},
\[
\mathsf D(S)
\]
is contained in every closed behaviour system containing \(S\). Hence it is
the smallest sub-Mealy system of \(\mathsf{Beh}\) containing \(S\).
\end{proof}

\section{Nerode quotient presentation}
\label{sec:nerode-quotient-presentation}

The residual closure has an effective quotient presentation. This
presentation identifies residual expressions precisely when they determine
the same behaviour.

\begin{definition}[Nerode equivalence]
\label{def:canonical-nerode-equivalence}
The \textbf{Nerode equivalence} on
\[
A_1\ltimes S
\]
is the kernel pair in
\[
\mathcal E/Z
\]
of the derivative evaluation map
\[
d_S:
A_1\ltimes S\to\mathsf{Beh}.
\]
It is denoted
\[
{\equiv_S}.
\]
Thus
\[
(a,\beta)\equiv_S(a',\beta')
\]
if and only if
\[
\partial_a\beta=\partial_{a'}\beta'
\]
as behaviours.
\end{definition}

\[
\begin{tikzcd}
{\equiv_S}
\arrow[r, shift left=0.5ex]
\arrow[r, shift right=0.5ex]
\arrow[d]
&
A_1\ltimes S
\arrow[d, "d_S"]
\\
A_1\ltimes S
\arrow[r, "d_S"']
&
\mathsf{Beh}
\arrow[ul, phantom, "\lrcorner", very near start]
\end{tikzcd}
\]

\begin{theorem}[Nerode quotient presentation]
\label{thm:canonical-nerode-quotient-presentation}
There is a canonical isomorphism in
\[
\mathcal E/Z
\]
\[
(A_1\ltimes S)/{\equiv_S}
\cong
\mathsf D(S).
\]
\end{theorem}

\begin{proof}
By definition,
\[
\mathsf D(S)
\]
is the image of
\[
d_S:
A_1\ltimes S\to\mathsf{Beh}.
\]
Factor \(d_S\) as a regular epimorphism followed by a monomorphism:
\[
A_1\ltimes S
\twoheadrightarrow
\mathsf D(S)
\hookrightarrow
\mathsf{Beh}.
\]
Because the second map is monic, the kernel pair of \(d_S\) is the kernel
pair of the regular epimorphism
\[
A_1\ltimes S
\twoheadrightarrow
\mathsf D(S).
\]
Since
\[
\mathcal E/Z
\]
is exact, this regular epimorphism is the effective quotient of its kernel
pair. Therefore
\[
(A_1\ltimes S)/{\equiv_S}
\cong
\mathsf D(S).
\]
\end{proof}

\[
\begin{tikzcd}[column sep=large]
{\equiv_S}
\arrow[r, shift left=0.5ex]
\arrow[r, shift right=0.5ex]
&
A_1\ltimes S
\arrow[r, two heads]
\arrow[rr, bend left=18, "d_S"]
&
(A_1\ltimes S)/{\equiv_S}
\arrow[r, hook, "\cong" description]
&
\mathsf{Beh}
\\
&&
\mathsf D(S)
\arrow[u, "\cong"']
\arrow[ur, hook']
&
\end{tikzcd}
\]

\begin{proposition}[Dependence only on image]
\label{prop:specifications-depend-only-on-image}
Let
\[
k:K\to\mathsf{Beh}
\qquad\text{and}\qquad
k':K'\to\mathsf{Beh}
\]
be morphisms in
\[
\mathcal E/Z
\]
with the same image subobject
\[
S\hookrightarrow\mathsf{Beh}.
\]
Then their canonical residual closures and Nerode quotient presentations are
canonically isomorphic to
\[
\mathsf D(S).
\]
\end{proposition}

\begin{proof}
The residual closure associated to a raw specification is the regular image of
all derivatives of behaviours in its image. Since both raw specifications
have the same image
\[
S\hookrightarrow\mathsf{Beh},
\]
both residual closures identify with
\[
\mathsf D(S).
\]

For the quotient presentation, the corresponding residual-expression objects
may differ, but each quotient is the effective quotient of its derivative
evaluation map. By
Theorem~\ref{thm:canonical-nerode-quotient-presentation}, each quotient is
canonically isomorphic to the image of derivative evaluation, namely
\[
\mathsf D(S).
\]
\end{proof}

\section{The residual-expression Mealy system}
\label{sec:residual-expression-mealy-system}

The object
\[
A_1\ltimes S
\]
does more than present the residual closure as an image. It carries a canonical
Mealy-system structure whose quotient by behavioural equality is
\[
\mathsf D(S).
\]

\begin{definition}[Residual-expression Mealy system]
\label{def:residual-expression-mealy-system}
Let
\[
S\hookrightarrow\mathsf{Beh}
\]
be a behaviour specification. The \textbf{residual-expression Mealy system}
\[
\mathsf E(S)
\]
has carrier
\[
A_1\ltimes S
\]
over \(Z\), with the slice structure of
Lemma~\ref{lem:residual-expression-slice-structure}. Its transition map is
defined by
\[
\delta_{\mathsf E(S)}(c,(a,\beta))
=
(a\circ c,\beta),
\]
where
\[
a:u\to v,
\qquad
c:w\to u,
\qquad
\beta\in S
\]
lies over \(v\). Its output map is
\[
\omega_{\mathsf E(S)}(c,(a,\beta))
=
\omega_{\mathsf{Beh}}(c,\partial_a\beta).
\]
\end{definition}

The typing is as follows. The state
\[
(a,\beta)
\]
has base
\[
\bigl(s_A(a),q_{\mathsf{Beh}}(\partial_a\beta)\bigr).
\]
An admissible input arrow
\[
c:w\to s_A(a)
\]
sends it to
\[
(a\circ c,\beta),
\]
whose behaviour is
\[
\partial_{a\circ c}\beta
=
\partial_c(\partial_a\beta).
\]
The output arrow is therefore
\[
q_{\mathsf{Beh}}(\partial_{a\circ c}\beta)
\to
q_{\mathsf{Beh}}(\partial_a\beta)
\]
in \(\mathbb B\), as required.

\begin{proposition}
\label{prop:residual-expression-is-mealy}
The data of Definition~\ref{def:residual-expression-mealy-system} define a
Mealy system
\[
\mathsf E(S):\mathbb A\rightsquigarrow\mathbb B.
\]
Moreover, the derivative evaluation map
\[
d_S:\mathsf E(S)\to\mathsf{Beh}
\]
is a Mealy transformation.
\end{proposition}

\begin{proof}
The transition unit law follows from
\[
a\circ e_A(s_Aa)=a.
\]
For composable arrows
\[
c:w\to u,
\qquad
h:z\to w,
\]
one has
\[
\delta_{\mathsf E(S)}(h,\delta_{\mathsf E(S)}(c,(a,\beta)))
=
\delta_{\mathsf E(S)}(h,(a\circ c,\beta))
=
(a\circ c\circ h,\beta),
\]
while
\[
\delta_{\mathsf E(S)}(c\circ h,(a,\beta))
=
(a\circ(c\circ h),\beta).
\]
Associativity in \(\mathbb A\) gives equality.

The output unit law follows from the output unit law in
\[
\mathsf{Beh}.
\]
For multiplication, the desired equation is
\[
\omega_{\mathsf E(S)}(c\circ h,(a,\beta))
=
\omega_{\mathsf E(S)}(c,(a,\beta))
\circ
\omega_{\mathsf E(S)}(h,(a\circ c,\beta)).
\]
Expanding the definitions, this is
\[
\omega_{\mathsf{Beh}}(c\circ h,\partial_a\beta)
=
\omega_{\mathsf{Beh}}(c,\partial_a\beta)
\circ
\omega_{\mathsf{Beh}}(h,\partial_c\partial_a\beta),
\]
which is exactly the output multiplication law in
\[
\mathsf{Beh}.
\]

Finally,
\[
d_S(\delta_{\mathsf E(S)}(c,(a,\beta)))
=
d_S(a\circ c,\beta)
=
\partial_{a\circ c}\beta
=
\partial_c(\partial_a\beta)
=
\delta_{\mathsf{Beh}}(c,d_S(a,\beta)).
\]
The output preservation equation is true by definition:
\[
\omega_{\mathsf E(S)}(c,(a,\beta))
=
\omega_{\mathsf{Beh}}(c,d_S(a,\beta)).
\]
Hence \(d_S\) is a Mealy transformation.
\end{proof}

\begin{corollary}
\label{cor:residual-expression-quotient-mealy}
The regular epimorphism
\[
e_S:\mathsf E(S)\twoheadrightarrow\mathsf D(S)
\]
in the image factorization of
\[
d_S
\]
is a Mealy transformation, and
\[
\mathsf D(S)
\]
is the effective Mealy quotient of
\[
\mathsf E(S)
\]
by the Nerode equivalence
\[
{\equiv_S}.
\]
\end{corollary}

\begin{proof}
By Proposition~\ref{prop:residual-expression-is-mealy}, the composite
\[
\mathsf E(S)
\xrightarrow{d_S}
\mathsf{Beh}
\]
is a Mealy transformation. Since
\[
d_S
=
j_S e_S
\]
with \(j_S\) monic, the Mealy transition and output maps of
\[
\mathsf D(S)
\]
are the unique restrictions making
\[
e_S
\]
a Mealy transformation. The quotient claim follows from
Theorem~\ref{thm:canonical-nerode-quotient-presentation}.
\end{proof}

\section{The syntactic Mealy category}
\label{sec:syntactic-mealy-category}

The residual closure is the canonical state object generated by the
specification. The corresponding syntactic algebraic object is the labelled
input-discrete internal category obtained from this Mealy system.

\begin{definition}[Syntactic Mealy category]
\label{def:syntactic-mealy-category}
Let
\[
S\hookrightarrow\mathsf{Beh}
\]
be a behaviour specification, and write
\[
D:=\mathsf D(S).
\]
The \textbf{syntactic Mealy category} of \(S\) is the action category
\[
\operatorname{Syn}_{\mathbb A,\mathbb B}(S)
:=
\int_{\mathbb A}D.
\]
Thus
\[
\operatorname{Syn}_{\mathbb A,\mathbb B}(S)_0
=
D,
\]
and
\[
\operatorname{Syn}_{\mathbb A,\mathbb B}(S)_1
=
A_1\times_{t_A,A_0,p_D}D.
\]
A generalized arrow is a pair
\[
(a,\beta)
\]
with
\[
a:u\to v,
\qquad
\beta\in D
\]
lying over \(v\). Its source and target are
\[
s(a,\beta)=\beta,
\qquad
t(a,\beta)=\partial_a\beta.
\]
\end{definition}

The identity at
\[
\beta\in D
\]
is
\[
(e_A(p_D\beta),\beta).
\]
If
\[
a:u\to v,
\qquad
c:w\to u,
\qquad
\beta\in D
\]
lies over \(v\), then the composite of
\[
(a,\beta):\beta\to\partial_a\beta
\]
and
\[
(c,\partial_a\beta):\partial_a\beta\to\partial_c\partial_a\beta
\]
is
\[
(a\circ c,\beta):\beta\to\partial_{a\circ c}\beta.
\]
The derivative law
\[
\partial_c(\partial_a\beta)=\partial_{a\circ c}\beta
\]
makes this well typed.

\begin{definition}[Canonical labellings]
\label{def:syntactic-mealy-category-labellings}
The syntactic Mealy category has an input labelling
\[
I_S:
\operatorname{Syn}_{\mathbb A,\mathbb B}(S)
\to
\mathbb A^{\mathrm{op}}
\]
whose object map is
\[
p_D:D\to A_0
\]
and whose arrow map is the projection
\[
A_1\times_{t_A,A_0,p_D}D\to A_1.
\]
It has an output labelling
\[
\Omega_S:
\operatorname{Syn}_{\mathbb A,\mathbb B}(S)
\to
\mathbb B^{\mathrm{op}}
\]
whose object map is
\[
q_D:D\to B_0
\]
and whose arrow map is
\[
\omega_D:
A_1\times_{t_A,A_0,p_D}D\to B_1,
\]
regarded as an arrow of
\[
\mathbb B^{\mathrm{op}}.
\]
\end{definition}

\begin{proposition}
\label{prop:syntactic-category-labelled-input-discrete}
The triple
\[
\left(
\operatorname{Syn}_{\mathbb A,\mathbb B}(S),
I_S,
\Omega_S
\right)
\]
is a labelled input-discrete internal category over
\[
\mathbb A^{\mathrm{op}}
\qquad\text{and}\qquad
\mathbb B^{\mathrm{op}}.
\]
Equivalently, it is the labelled input-discrete category associated to the
Mealy system
\[
\mathsf D(S):\mathbb A\rightsquigarrow\mathbb B.
\]
\end{proposition}

\begin{proof}
The action category
\[
\int_{\mathbb A}D
\]
is an internal category by the action laws of the Mealy transition on \(D\).
The map
\[
I_S
\]
is an internal functor to
\[
\mathbb A^{\mathrm{op}},
\]
because an arrow
\[
(a,\beta):\beta\to\partial_a\beta
\]
lies over the arrow
\[
a:p_D(\partial_a\beta)\to p_D(\beta)
\]
of \(\mathbb A\), hence over an arrow
\[
p_D(\beta)\to p_D(\partial_a\beta)
\]
of
\[
\mathbb A^{\mathrm{op}}.
\]
It is input-discrete because the arrow object of
\[
\int_{\mathbb A}D
\]
is exactly the pullback
\[
A_1\times_{t_A,A_0,p_D}D.
\]

The output labelling
\[
\Omega_S
\]
is an internal functor to
\[
\mathbb B^{\mathrm{op}}
\]
because the Mealy output equations for \(D\) are precisely the identity and
composition preservation equations for this labelling. Explicitly,
\[
\omega_D(a,\beta):
q_D(\partial_a\beta)\to q_D(\beta)
\]
is an arrow of \(\mathbb B\), hence an arrow
\[
q_D(\beta)\to q_D(\partial_a\beta)
\]
of
\[
\mathbb B^{\mathrm{op}}.
\]
The output unit and multiplication laws are exactly functoriality.
\end{proof}

\section{Residual-expression categories and syntactic quotients}
\label{sec:residual-expression-categories-syntactic-quotients}

The residual-expression Mealy system
\[
\mathsf E(S)
\]
has its own action category. This is the internal category of formal residual
expressions before quotienting by behavioural equality.

\begin{definition}[Residual-expression category]
\label{def:residual-expression-category}
The \textbf{residual-expression category} of \(S\) is
\[
\operatorname{ResExpr}_{\mathbb A,\mathbb B}(S)
:=
\int_{\mathbb A}\mathsf E(S).
\]
Its object object is
\[
A_1\ltimes S.
\]
Its arrow object is
\[
A_1\times_{t_A,A_0,p_{A_1\ltimes S}}(A_1\ltimes S),
\]
where
\[
p_{A_1\ltimes S}=s_A\pi_1
\]
is the \(A_0\)-component induced by derivative evaluation.
A generalized arrow is a triple
\[
(c,a,\beta)
\]
with
\[
a:u\to v,
\qquad
c:w\to u,
\qquad
\beta\in S
\]
lying over \(v\). Its source and target are
\[
s(c,a,\beta)=(a,\beta),
\qquad
t(c,a,\beta)=(a\circ c,\beta).
\]
\end{definition}

The residual-expression category carries input and output labellings
\[
I_{\mathsf E(S)}:
\operatorname{ResExpr}_{\mathbb A,\mathbb B}(S)
\to
\mathbb A^{\mathrm{op}},
\]
and
\[
\Omega_{\mathsf E(S)}:
\operatorname{ResExpr}_{\mathbb A,\mathbb B}(S)
\to
\mathbb B^{\mathrm{op}}
\]
induced by the Mealy system
\[
\mathsf E(S).
\]
The output labelling sends
\[
(c,a,\beta)
\]
to
\[
\omega_{\mathsf{Beh}}(c,\partial_a\beta).
\]

\begin{proposition}[Quotient functor to the syntactic category]
\label{prop:residual-expression-category-quotient}
The Mealy quotient
\[
e_S:\mathsf E(S)\twoheadrightarrow\mathsf D(S)
\]
induces a labelled internal functor
\[
\int_{\mathbb A}e_S:
\operatorname{ResExpr}_{\mathbb A,\mathbb B}(S)
\to
\operatorname{Syn}_{\mathbb A,\mathbb B}(S).
\]
It is a regular epimorphism on object objects and on arrow objects. Its object
map is
\[
(a,\beta)\longmapsto \partial_a\beta,
\]
and its arrow map is
\[
(c,a,\beta)
\longmapsto
(c,\partial_a\beta).
\]
\end{proposition}

\begin{proof}
Because
\[
e_S:\mathsf E(S)\to\mathsf D(S)
\]
is a Mealy transformation, it induces an internal functor between the
associated action categories. The object map is the carrier map
\[
e_S:A_1\ltimes S\to\mathsf D(S),
\]
hence sends
\[
(a,\beta)
\]
to
\[
\partial_a\beta.
\]
The arrow map is obtained by pulling this map back along
\[
t_A:A_1\to A_0.
\]
Thus it sends
\[
(c,a,\beta)
\]
to
\[
(c,\partial_a\beta).
\]

The object map is a regular epimorphism by definition of the image
factorization. The arrow map is a pullback of this regular epimorphism along
\[
A_1\to A_0
\]
and is therefore again a regular epimorphism. Preservation of input and output
labellings follows from preservation of the input action and output map by
the Mealy transformation \(e_S\).
\end{proof}

\begin{definition}[Syntactic congruence on residual expressions]
\label{def:syntactic-congruence-residual-expression-category}
The \textbf{syntactic congruence} on
\[
\operatorname{ResExpr}_{\mathbb A,\mathbb B}(S)
\]
is the levelwise kernel pair of the labelled functor
\[
\int_{\mathbb A}e_S:
\operatorname{ResExpr}_{\mathbb A,\mathbb B}(S)
\to
\operatorname{Syn}_{\mathbb A,\mathbb B}(S).
\]
On objects it is the Nerode equivalence
\[
(a,\beta)\equiv_S(a',\beta')
\quad\Longleftrightarrow\quad
\partial_a\beta=\partial_{a'}\beta'.
\]
On arrows it identifies
\[
(c,a,\beta)
\]
and
\[
(c',a',\beta')
\]
precisely when
\[
c=c'
\]
as arrows of \(\mathbb A\), and
\[
\partial_a\beta=\partial_{a'}\beta'.
\]
\end{definition}

\begin{theorem}[Syntactic category as effective quotient]
\label{thm:syntactic-category-effective-quotient}
The syntactic Mealy category
\[
\operatorname{Syn}_{\mathbb A,\mathbb B}(S)
\]
is the effective quotient of the residual-expression category
\[
\operatorname{ResExpr}_{\mathbb A,\mathbb B}(S)
\]
by its syntactic congruence. This quotient preserves the input and output
labellings, so it is an effective quotient in labelled input-discrete internal
categories.
\end{theorem}

\begin{proof}
The object quotient is
\[
A_1\ltimes S\twoheadrightarrow\mathsf D(S),
\]
which is the effective quotient of the object-level Nerode equivalence by
Theorem~\ref{thm:canonical-nerode-quotient-presentation}.

The arrow map of
\[
\int_{\mathbb A}e_S
\]
is obtained by pulling back this regular epimorphism along the fixed map
\[
t_A:A_1\to A_0.
\]
Since regular epimorphisms and their effective kernel-pair quotients are
stable under pullback in a topos, the arrow map is the effective quotient of
its kernel pair. Hence the quotient is effective levelwise.

The source, target, identity, and composition maps of
\[
\operatorname{Syn}_{\mathbb A,\mathbb B}(S)
\]
are the unique maps induced from those of
\[
\operatorname{ResExpr}_{\mathbb A,\mathbb B}(S)
\]
through the effective quotients. These maps are well-defined because
\[
e_S
\]
is a Mealy transformation. The input and output labellings also descend
because
\[
\int_{\mathbb A}e_S
\]
preserves them. Therefore the quotient is an effective quotient in labelled
input-discrete internal categories.
\end{proof}

\begin{remark}[Levelwise effectiveness]
\label{rem:levelwise-effectiveness-syntactic-category}
Effectiveness in
Theorem~\ref{thm:syntactic-category-effective-quotient} is understood
levelwise on object and arrow objects. The internal-category structure and
the input and output labellings are induced by exactness and by the fact that
the quotient functor preserves the Mealy transition and output data.
\end{remark}

\section{Generated recognizers and the universal property}
\label{sec:generated-recognizers-universal-property}

The syntactic Mealy category is canonical because every generated recognizer
of the specification maps onto it, and every generated separated recognizer is
isomorphic to it.

\begin{definition}[Pointed realization]
\label{def:pointed-realization}
A \textbf{pointed realization} of a behaviour specification
\[
S\hookrightarrow\mathsf{Beh}
\]
is a Mealy system \(X\), together with a morphism
\[
i:S\to X
\]
in
\[
\mathcal E/Z
\]
such that the composite
\[
S\xrightarrow{i}X\xrightarrow{\beh_X}\mathsf{Beh}
\]
is the inclusion
\[
S\hookrightarrow\mathsf{Beh}.
\]
\end{definition}

\[
\begin{tikzcd}[column sep=large]
S
\arrow[r, "i"]
\arrow[dr, hook']
&
X
\arrow[d, "\beh_X"]
\\
&
\mathsf{Beh}
\end{tikzcd}
\]

Thus a pointed realization includes a chosen realizing state of \(X\) for each
specified behaviour.

\begin{definition}[Generated part]
\label{def:reachable-part-generated-by-pointed-specification}
Let
\[
i:S\to X
\]
be a pointed realization. Define the \textbf{generated part} of \(X\) from
\(S\) to be the image in
\[
\mathcal E/Z
\]
of the map
\[
\rho_{X,S}:
A_1\ltimes S\to X
\]
given by
\[
\rho_{X,S}(a,\beta)=\delta_X(a,i(\beta)).
\]
It is denoted
\[
\operatorname{Gen}_X(S)\hookrightarrow X.
\]
\end{definition}

\[
\begin{tikzcd}[column sep=large]
A_1\ltimes S
\arrow[rr, "\rho_{X,S}"]
\arrow[dr, two heads]
&&
X
\\
&
\operatorname{Gen}_X(S)
\arrow[ur, hook']
&
\end{tikzcd}
\]

\begin{lemma}[Typing of the generated map]
\label{lem:generated-map-slice-typing}
The map
\[
\rho_{X,S}:A_1\ltimes S\to X
\]
is a morphism in
\[
\mathcal E/Z,
\]
where \(A_1\ltimes S\) is given the slice structure induced by
\[
d_S.
\]
\end{lemma}

\begin{proof}
The behaviour map
\[
\beh_X:X\to\mathsf{Beh}
\]
is a Mealy transformation. Hence it preserves derivatives:
\[
\beh_X(\delta_X(a,i(\beta)))
=
\partial_a(\beh_X(i(\beta))).
\]
Since \(i\) is a pointed realization,
\[
\beh_X(i(\beta))=\beta.
\]
Therefore
\[
\beh_X(\rho_{X,S}(a,\beta))
=
\partial_a\beta
=
d_S(a,\beta).
\]
Because \(\beh_X\) is a morphism over \(Z\), this says exactly that the base
point of
\[
\rho_{X,S}(a,\beta)
\]
is the base point of
\[
d_S(a,\beta).
\]
Thus \(\rho_{X,S}\) is a morphism in
\[
\mathcal E/Z.
\]
\end{proof}

\[
\begin{tikzcd}[column sep=large]
A_1\ltimes S
\arrow[r, "\rho_{X,S}"]
\arrow[d, "d_S"']
&
X
\arrow[d, "\beh_X"]
\\
\mathsf{Beh}
\arrow[r, equal]
&
\mathsf{Beh}
\end{tikzcd}
\]

\begin{proposition}[Generated part is a sub-Mealy system]
\label{prop:generated-part-submealy}
Let
\[
i:S\to X
\]
be a pointed realization. Then
\[
\operatorname{Gen}_X(S)\hookrightarrow X
\]
is a sub-Mealy system of \(X\).
\end{proposition}

\begin{proof}
Let
\[
A_1\ltimes S
\twoheadrightarrow
\operatorname{Gen}_X(S)
\hookrightarrow
X
\]
be the image factorization of
\[
\rho_{X,S}.
\]
Since regular epimorphisms are stable under pullback, the induced map
\[
A_1\times_{A_0}(A_1\ltimes S)
\twoheadrightarrow
A_1\times_{A_0}\operatorname{Gen}_X(S)
\]
is a regular epimorphism, where both pullbacks use the \(A_0\)-components
induced by the relevant slice structures.

A generalized element of the upper-left object is a triple
\[
(c,a,\beta)
\]
with
\[
a:u\to v,
\qquad
c:w\to u,
\qquad
\beta\in S
\]
lying over \(v\). After pulling back along the displayed regular epimorphism,
the transition of \(X\) sends this element to
\[
\delta_X(c,\delta_X(a,i(\beta))).
\]
By the Mealy transition law, this equals
\[
\delta_X(a\circ c,i(\beta)).
\]
This lies in the image of \(\rho_{X,S}\), via the residual expression
\[
(a\circ c,\beta).
\]
Therefore the transition of \(X\), after pullback along a regular epimorphism,
lands in
\[
\operatorname{Gen}_X(S).
\]
Pullback along a regular epimorphism reflects containment of subobjects in a
regular category, so the transition of \(X\) restricts to
\[
A_1\times_{A_0}\operatorname{Gen}_X(S)
\to
\operatorname{Gen}_X(S).
\]

The output map restricts from the output map of \(X\), and the Mealy equations
hold because they hold in \(X\). Hence
\[
\operatorname{Gen}_X(S)
\]
is a sub-Mealy system of \(X\).
\end{proof}

\begin{proposition}[Generated part maps onto residual closure]
\label{prop:generated-part-to-residual-closure}
There is a canonical regular epimorphism
\[
\operatorname{Gen}_X(S)
\twoheadrightarrow
\mathsf D(S).
\]
\end{proposition}

\begin{proof}
By Lemma~\ref{lem:generated-map-slice-typing}, the composite
\[
A_1\ltimes S
\xrightarrow{\rho_{X,S}}
X
\xrightarrow{\beh_X}
\mathsf{Beh}
\]
is precisely
\[
d_S:
A_1\ltimes S\to\mathsf{Beh}.
\]
Factor
\[
\rho_{X,S}
\]
through its image:
\[
A_1\ltimes S
\twoheadrightarrow
\operatorname{Gen}_X(S)
\hookrightarrow
X.
\]
The image of the composite
\[
A_1\ltimes S\to X\to\mathsf{Beh}
\]
is therefore the image of the induced map
\[
\operatorname{Gen}_X(S)\to\mathsf{Beh}.
\]
But this composite is \(d_S\), whose image is
\[
\mathsf D(S).
\]
Hence the induced map
\[
\operatorname{Gen}_X(S)\to\mathsf D(S)
\]
is a regular epimorphism.
\end{proof}

\[
\begin{tikzcd}[column sep=large, row sep=large]
A_1\ltimes S
\arrow[r, "\rho_{X,S}"]
\arrow[d, "d_S"']
\arrow[dr, two heads, "\pi_X" description]
&
X
\arrow[d, "\beh_X"]
\\
\mathsf{Beh}
&
\operatorname{Gen}_X(S)
\arrow[u, hook']
\arrow[l, dashed, two heads, "\overline{\beh}_X"']
\\
\mathsf D(S)
\arrow[u, hook]
&
\end{tikzcd}
\]

\begin{definition}[Generated and separated pointed realization]
\label{def:generated-separated-pointed-realization}
A pointed realization
\[
i:S\to X
\]
is \textbf{generated by \(S\)} if the inclusion
\[
\operatorname{Gen}_X(S)\hookrightarrow X
\]
is an isomorphism.

It is \textbf{behaviourally separated} if the behaviour map
\[
\beh_X:X\to\mathsf{Beh}
\]
is monic.
\end{definition}

\begin{theorem}[Generated separated realizations are syntactic]
\label{thm:generated-separated-realization-residual-closure}
If a pointed realization \(X\) is generated by \(S\) and behaviourally
separated, then
\[
X\cong \mathsf D(S)
\]
as Mealy systems. Consequently
\[
\int_{\mathbb A}X
\cong
\operatorname{Syn}_{\mathbb A,\mathbb B}(S)
\]
as labelled input-discrete internal categories.
\end{theorem}

\begin{proof}
By Proposition~\ref{prop:generated-part-to-residual-closure}, there is a
regular epimorphism
\[
\operatorname{Gen}_X(S)
\twoheadrightarrow
\mathsf D(S).
\]
If \(X\) is generated by \(S\), then
\[
\operatorname{Gen}_X(S)\cong X,
\]
so there is a regular epimorphism
\[
X\twoheadrightarrow\mathsf D(S).
\]
The composite
\[
X\twoheadrightarrow\mathsf D(S)\hookrightarrow\mathsf{Beh}
\]
is the behaviour map of \(X\). If this behaviour map is monic, then the
regular epimorphism
\[
X\twoheadrightarrow\mathsf D(S)
\]
is also monic. A morphism that is both regular epi and mono is an isomorphism.

The morphism
\[
X\to\mathsf D(S)
\]
is a Mealy transformation because it is induced by the behaviour map. Its
inverse is also a Mealy transformation because
\[
\mathsf D(S)
\]
has the restricted Mealy structure inherited from the final behaviour system.
Thus
\[
X\cong\mathsf D(S)
\]
as Mealy systems.

Applying the action-category construction
\[
\int_{\mathbb A}(-)
\]
gives an isomorphism of the corresponding labelled input-discrete internal
categories:
\[
\int_{\mathbb A}X
\cong
\int_{\mathbb A}\mathsf D(S)
=
\operatorname{Syn}_{\mathbb A,\mathbb B}(S).
\]
\end{proof}

\begin{theorem}[Syntactic quotient of generated recognizers]
\label{thm:syntactic-quotient-generated-recognizers}
Let
\[
i:S\to X
\]
be a pointed realization. Then the canonical regular epimorphism
\[
\operatorname{Gen}_X(S)\twoheadrightarrow\mathsf D(S)
\]
induces a labelled internal functor
\[
\int_{\mathbb A}\operatorname{Gen}_X(S)
\twoheadrightarrow
\operatorname{Syn}_{\mathbb A,\mathbb B}(S)
\]
which is a regular epimorphism on objects and arrows. Hence the syntactic
Mealy category is a quotient of the labelled input-discrete category generated
inside any pointed realization of \(S\).
\end{theorem}

\begin{proof}
By Proposition~\ref{prop:generated-part-submealy},
\[
\operatorname{Gen}_X(S)
\]
is a sub-Mealy system of \(X\). The regular epimorphism
\[
\operatorname{Gen}_X(S)\twoheadrightarrow\mathsf D(S)
\]
is induced by the behaviour map of \(X\), hence is a Mealy transformation. It
therefore induces a labelled internal functor between action categories.

The object map is a regular epimorphism by
Proposition~\ref{prop:generated-part-to-residual-closure}. The arrow map is
its pullback along
\[
t_A:A_1\to A_0,
\]
and is therefore again a regular epimorphism. The induced functor preserves
input and output labellings because the original map is a Mealy
transformation.
\end{proof}

\section{The one-object monoid case}
\label{sec:canonical-residual-closure-one-object}

Let
\[
\mathcal E=\mathbf{Set},
\]
and let
\[
\mathbb A=M,
\qquad
\mathbb B=N
\]
be one-object internal categories, hence monoids.

We keep the convention
\[
mn=m\circ n.
\]
Thus \(mn\) means first \(n\), then \(m\) as categorical composition, while
the right-action execution by \(mn\) is computed as first \(m\), then \(n\).

A behaviour is an \(N\)-valued residual cocycle
\[
\lambda:M\times M\to N
\]
satisfying
\[
\lambda(r,1)=1,
\qquad
\lambda(r,mn)=\lambda(r,m)\lambda(rm,n).
\]
The derivative by
\[
r\in M
\]
is
\[
(\partial_r\lambda)(s,m)=\lambda(rs,m).
\]

\begin{proposition}[Residual closure in the monoid case]
\label{prop:one-object-residual-closure}
Let
\[
S\subseteq \mathsf{Beh}_{M,N}
\]
be a behaviour specification. Then
\[
\mathsf D(S)
=
\{\partial_r\lambda\mid r\in M,\ \lambda\in S\}.
\]
\end{proposition}

\begin{proof}
The derivative evaluation map is
\[
M\times S\to \mathsf{Beh}_{M,N},
\qquad
(r,\lambda)\longmapsto \partial_r\lambda.
\]
Its image is exactly the displayed set.
\end{proof}

\begin{corollary}[Singleton residual closure]
\label{cor:one-object-singleton-residual-closure}
For a single behaviour
\[
\lambda\in\mathsf{Beh}_{M,N},
\]
one has
\[
\mathsf D(\lambda)
:=
\mathsf D(\{\lambda\})
=
\{\partial_r\lambda\mid r\in M\}.
\]
\end{corollary}

\begin{corollary}[Nerode quotient in the monoid case]
\label{cor:one-object-nerode-quotient}
For a single behaviour
\[
\lambda\in\mathsf{Beh}_{M,N},
\]
define
\[
r\equiv_\lambda s
\]
if and only if
\[
\partial_r\lambda=\partial_s\lambda.
\]
Equivalently,
\[
r\equiv_\lambda s
\]
if and only if, for all
\[
u,m\in M,
\]
one has
\[
\lambda(ru,m)=\lambda(su,m).
\]
Then
\[
M/{\equiv_\lambda}
\cong
\mathsf D(\lambda).
\]
\end{corollary}

\begin{proof}
For a singleton specification, the residual-expression object is \(M\), and
the derivative evaluation map is
\[
M\to\mathsf{Beh}_{M,N},
\qquad
r\longmapsto \partial_r\lambda.
\]
Its kernel pair is exactly
\[
{\equiv_\lambda},
\]
and its image is
\[
\mathsf D(\lambda).
\]
The result follows from
Theorem~\ref{thm:canonical-nerode-quotient-presentation}.
\end{proof}

\section{The syntactic transition-output monoid}
\label{sec:syntactic-transition-output-monoid}

In the one-object case, the syntactic Mealy category has an associated
transition-output monoid. This is the algebraic object that records how each
input monoid element acts on the canonical residual state set, together with
the output it produces at each residual state.

Let
\[
S\subseteq\mathsf{Beh}_{M,N}
\]
and write
\[
D:=\mathsf D(S).
\]
The Mealy system \(D\) is a right \(M\)-set with output cocycle
\[
\omega_D:M\times D\to N.
\]
For each
\[
m\in M,
\]
define
\[
\tau_m:D\to D,
\qquad
\tau_m(\lambda)=\partial_m\lambda,
\]
and
\[
o_m:D\to N,
\qquad
o_m(\lambda)=\omega_D(m,\lambda).
\]
If
\[
\lambda=\partial_r\lambda_0
\]
with
\[
\lambda_0\in S,
\]
then
\[
o_m(\partial_r\lambda_0)=\lambda_0(r,m).
\]

\begin{definition}[Mealy transformation monoid]
\label{def:mealy-transformation-monoid}
Let
\[
\operatorname{MealyEnd}_N(D)
\]
be the set of pairs
\[
(o,\tau)
\]
where
\[
\tau:D\to D
\qquad\text{and}\qquad
o:D\to N.
\]
The multiplication is
\[
(o,\tau)\cdot(o',\tau')
=
\bigl(\lambda\mapsto o(\lambda)o'(\tau\lambda),\ \tau'\circ\tau\bigr),
\]
and the unit is
\[
(\lambda\mapsto 1_N,\operatorname{id}_D).
\]
\end{definition}

\begin{remark}[Operational order]
\label{rem:mealy-end-operational-order}
The multiplication in
\[
\operatorname{MealyEnd}_N(D)
\]
is written in the same operational order as the right \(M\)-action. The
product
\[
(o,\tau)\cdot(o',\tau')
\]
means first perform \((o,\tau)\), then perform \((o',\tau')\). Hence the state
component is
\[
\tau'\circ\tau,
\]
while the output component first emits
\[
o(\lambda)
\]
and then emits
\[
o'(\tau\lambda).
\]
\end{remark}

\begin{proposition}
\label{prop:mealy-transformation-monoid}
The multiplication in
\[
\operatorname{MealyEnd}_N(D)
\]
is associative and unital. Hence
\[
\operatorname{MealyEnd}_N(D)
\]
is a monoid.
\end{proposition}

\begin{proof}
Let
\[
(o,\tau),
\qquad
(o',\tau'),
\qquad
(o'',\tau'')
\]
be three elements. The transition components associate because composition of
functions is associative:
\[
\tau''\circ(\tau'\circ\tau)
=
(\tau''\circ\tau')\circ\tau.
\]
The output component of
\[
\bigl((o,\tau)\cdot(o',\tau')\bigr)\cdot(o'',\tau'')
\]
at \(\lambda\) is
\[
o(\lambda)o'(\tau\lambda)o''(\tau'\tau\lambda).
\]
The output component of
\[
(o,\tau)\cdot\bigl((o',\tau')\cdot(o'',\tau'')\bigr)
\]
at \(\lambda\) is the same expression. Associativity in \(N\) gives equality.
The stated unit is immediate from the unit of \(N\) and the identity function
on \(D\).
\end{proof}

\begin{proposition}[Input representation]
\label{prop:input-representation-transition-output}
The assignment
\[
\chi_S:M\to\operatorname{MealyEnd}_N(D),
\qquad
m\longmapsto(o_m,\tau_m),
\]
is a monoid homomorphism.
\end{proposition}

\begin{proof}
For
\[
m,n\in M
\]
and
\[
\lambda\in D,
\]
the right-action law gives
\[
\tau_{mn}(\lambda)
=
\partial_{mn}\lambda
=
\partial_n(\partial_m\lambda)
=
\tau_n(\tau_m\lambda).
\]
Thus
\[
\tau_{mn}
=
\tau_n\circ\tau_m.
\]
The output cocycle law gives
\[
o_{mn}(\lambda)
=
\omega_D(mn,\lambda)
=
\omega_D(m,\lambda)\omega_D(n,\lambda\cdot m)
=
o_m(\lambda)o_n(\tau_m\lambda).
\]
Therefore
\[
(o_{mn},\tau_{mn})
=
(o_m,\tau_m)\cdot(o_n,\tau_n),
\]
so \(\chi_S\) is a monoid homomorphism.
\end{proof}

\begin{definition}[Syntactic transition-output monoid]
\label{def:syntactic-transition-output-monoid}
The \textbf{syntactic transition-output monoid} of
\[
S\subseteq\mathsf{Beh}_{M,N}
\]
is the image
\[
\operatorname{SynMon}_{M,N}(S)
:=
\operatorname{im}
\left(
M
\xrightarrow{\chi_S}
\operatorname{MealyEnd}_N(\mathsf D(S))
\right).
\]
\end{definition}

\begin{definition}[Syntactic congruence on \(M\)]
\label{def:syntactic-congruence-on-input-monoid}
Define a congruence
\[
{\equiv_S^{\operatorname{syn}}}
\]
on \(M\) by
\[
m\equiv_S^{\operatorname{syn}} n
\]
if and only if
\[
(o_m,\tau_m)=(o_n,\tau_n)
\]
as elements of
\[
\operatorname{MealyEnd}_N(\mathsf D(S)).
\]
Equivalently,
\[
m\equiv_S^{\operatorname{syn}} n
\]
if and only if, for every
\[
\lambda_0\in S
\]
and every
\[
r\in M,
\]
one has
\[
\partial_{rm}\lambda_0
=
\partial_{rn}\lambda_0
\]
and
\[
\lambda_0(r,m)=\lambda_0(r,n).
\]
\end{definition}

\begin{theorem}[Syntactic monoid quotient]
\label{thm:syntactic-monoid-quotient}
There is a canonical monoid isomorphism
\[
M/{\equiv_S^{\operatorname{syn}}}
\cong
\operatorname{SynMon}_{M,N}(S).
\]
\end{theorem}

\begin{proof}
The congruence
\[
{\equiv_S^{\operatorname{syn}}}
\]
is by definition the kernel congruence of the monoid homomorphism
\[
\chi_S:
M\to
\operatorname{MealyEnd}_N(\mathsf D(S)).
\]
The image of this homomorphism is
\[
\operatorname{SynMon}_{M,N}(S).
\]
Therefore the usual image-kernel factorization of monoids gives
\[
M/{\equiv_S^{\operatorname{syn}}}
\cong
\operatorname{SynMon}_{M,N}(S).
\]
\end{proof}

\begin{remark}[State Nerode versus syntactic congruence]
\label{rem:state-nerode-versus-syntactic-congruence}
For a singleton behaviour
\[
\lambda,
\]
the Nerode quotient
\[
M/{\equiv_\lambda}
\]
identifies residual histories that determine the same residual behaviour:
\[
r\equiv_\lambda s
\quad\Longleftrightarrow\quad
\partial_r\lambda=\partial_s\lambda.
\]
The syntactic congruence
\[
{\equiv_\lambda^{\operatorname{syn}}}
\]
instead identifies input elements that induce the same transition-output
operation on all residual states. Thus it remembers both the derivative it
produces and the output it emits at each residual state. This is the
Mealy-output refinement of the usual transition congruence of the minimal
automaton.
\end{remark}

\section{Free-input normal form}
\label{sec:canonical-residual-closure-free-input}

Suppose now that
\[
M=I^\ast
\]
is the free monoid on a set \(I\). Then the final behaviour object has the
classical generator-level normal form
\[
\mathsf{Beh}_{I^\ast,N}
\cong
N^{I^+},
\]
where
\[
I^+
\]
is the set of nonempty words.

Under this isomorphism, a behaviour is a function
\[
\beta:I^+\to N.
\]
The derivative by a word
\[
w\in I^\ast
\]
is the residual shift
\[
(\partial_w\beta)(u)=\beta(wu),
\qquad
u\in I^+.
\]

\begin{proposition}[Residual closure in free-input normal form]
\label{prop:free-input-residual-closure}
For
\[
S\subseteq N^{I^+},
\]
one has
\[
\mathsf D(S)
=
\{\partial_w\beta\mid w\in I^\ast,\ \beta\in S\}.
\]
In particular, for a single behaviour
\[
\beta:I^+\to N,
\]
one has
\[
\mathsf D(\beta)
=
\{\partial_w\beta\mid w\in I^\ast\}.
\]
\end{proposition}

\begin{proof}
This is the specialization of
Proposition~\ref{prop:one-object-residual-closure} under the isomorphism
\[
\mathsf{Beh}_{I^\ast,N}\cong N^{I^+}.
\]
The derivative operation becomes the residual shift
\[
(\partial_w\beta)(u)=\beta(wu).
\]
\end{proof}

For a behaviour
\[
\beta:I^+\to N,
\]
write
\[
\lambda_\beta:I^\ast\times I^\ast\to N
\]
for the corresponding coherent residual cocycle. Thus
\[
\lambda_\beta(r,\epsilon)=1,
\]
and if
\[
u=i_1\cdots i_n,
\qquad n\geq 1,
\]
then
\[
\lambda_\beta(r,u)
=
\beta(ri_1)\beta(ri_1i_2)\cdots\beta(ri_1\cdots i_n).
\]

\begin{theorem}[Free-input syntactic transition-output monoid]
\label{thm:free-input-syntactic-transition-output-monoid}
Let
\[
S\subseteq N^{I^+}.
\]
The syntactic transition-output monoid of \(S\) is the image of
\[
I^\ast
\to
\operatorname{MealyEnd}_N(\mathsf D(S)),
\]
which sends a word
\[
u\in I^\ast
\]
to the pair
\[
(o_u,\tau_u),
\]
where
\[
\tau_u(\beta)=\partial_u\beta
\]
and
\[
o_u(\partial_r\beta)=\lambda_\beta(r,u).
\]
Equivalently, it is the quotient of
\[
I^\ast
\]
by the congruence
\[
u\equiv_S^{\operatorname{syn}} v
\]
if and only if, for every
\[
\beta\in S
\]
and every
\[
r\in I^\ast,
\]
one has
\[
\partial_{ru}\beta=\partial_{rv}\beta
\]
and
\[
\lambda_\beta(r,u)=\lambda_\beta(r,v).
\]
\end{theorem}

\begin{proof}
This is Theorem~\ref{thm:syntactic-monoid-quotient} specialized to the free
monoid
\[
I^\ast.
\]
The only additional point is the expression for the output component. Under
the isomorphism
\[
\mathsf{Beh}_{I^\ast,N}\cong N^{I^+},
\]
the output produced by processing a word
\[
u
\]
after residual history
\[
r
\]
is the coherent cocycle value
\[
\lambda_\beta(r,u).
\]
Thus equality of output functions is exactly the displayed condition
\[
\lambda_\beta(r,u)=\lambda_\beta(r,v)
\]
for all residual prefixes \(r\) and all specified behaviours \(\beta\).
\end{proof}

\begin{remark}[Classical derivative automata]
\label{rem:free-input-classical-derivative-automata}
When the output monoid is specialized to the usual observation structure of
deterministic language recognition, the residual closure
\[
\mathsf D(\beta)
=
\{\partial_w\beta\mid w\in I^\ast\}
\]
is the state set of the minimal derivative automaton. The syntactic
transition-output monoid above is the transition monoid of that minimal
residual automaton, with the Mealy output component retained rather than
discarded.
\end{remark}

\section{Main consolidation}
\label{sec:canonical-residual-closure-main-consolidation}

\begin{theorem}[Canonical residual closure and syntactic Mealy category]
\label{thm:canonical-residual-closure-main}
Let
\[
\mathcal E
\]
be a Grothendieck topos, and let
\[
\mathbb A,\mathbb B
\]
be internal categories in \(\mathcal E\). Let
\[
\mathsf{Beh}_{\mathbb A,\mathbb B}
\]
be the final behaviour system for Mealy morphisms
\[
\mathbb A\rightsquigarrow\mathbb B.
\]
Then:
\begin{enumerate}
\item Every raw behaviour specification
\[
k:K\to\mathsf{Beh}_{\mathbb A,\mathbb B}
\]
has an image specification
\[
\operatorname{im}(k)\hookrightarrow\mathsf{Beh}_{\mathbb A,\mathbb B}
\]
in
\[
\mathcal E/(A_0\times B_0).
\]

\item For every behaviour subobject
\[
S\hookrightarrow\mathsf{Beh}_{\mathbb A,\mathbb B},
\]
there is a derivative evaluation map
\[
d_S:
A_1\ltimes S\to\mathsf{Beh}_{\mathbb A,\mathbb B}.
\]

\item The residual closure
\[
\mathsf D(S)
=
\operatorname{im}
\left(
A_1\ltimes S
\xrightarrow{d_S}
\mathsf{Beh}_{\mathbb A,\mathbb B}
\right)
\]
defines a closure operator on
\[
\operatorname{Sub}_{A_0\times B_0}
(\mathsf{Beh}_{\mathbb A,\mathbb B}).
\]

\item The closed subobjects for this closure operator are precisely the
sub-Mealy systems of the final behaviour system.

\item The inclusion
\[
\operatorname{ClosedBeh}_{\mathbb A,\mathbb B}
\hookrightarrow
\operatorname{Sub}_{A_0\times B_0}
(\mathsf{Beh}_{\mathbb A,\mathbb B})
\]
has left adjoint
\[
\mathsf D.
\]
Equivalently, for every closed behaviour system \(D\),
\[
\mathsf D(S)\leq D
\quad\Longleftrightarrow\quad
S\leq D.
\]

\item The canonical residual realization has an effective quotient
presentation
\[
\mathsf D(S)
\cong
(A_1\ltimes S)/{\equiv_S},
\]
where
\[
(a,\beta)\equiv_S(a',\beta')
\]
if and only if
\[
\partial_a\beta=\partial_{a'}\beta'.
\]

\item The residual-expression object
\[
A_1\ltimes S
\]
carries a canonical residual-expression Mealy system
\[
\mathsf E(S),
\]
and the quotient map
\[
\mathsf E(S)\twoheadrightarrow\mathsf D(S)
\]
is a Mealy transformation and an effective quotient by behavioural equality.

\item The syntactic Mealy category of \(S\) is
\[
\operatorname{Syn}_{\mathbb A,\mathbb B}(S)
=
\int_{\mathbb A}\mathsf D(S),
\]
equipped with canonical labellings
\[
I_S:
\operatorname{Syn}_{\mathbb A,\mathbb B}(S)\to\mathbb A^{\mathrm{op}},
\qquad
\Omega_S:
\operatorname{Syn}_{\mathbb A,\mathbb B}(S)\to\mathbb B^{\mathrm{op}}.
\]

\item The syntactic Mealy category is the effective quotient of the
residual-expression category
\[
\operatorname{ResExpr}_{\mathbb A,\mathbb B}(S)
=
\int_{\mathbb A}\mathsf E(S)
\]
by the induced syntactic congruence.

\item If \(X\) is a pointed realization of \(S\), then its generated part is a
sub-Mealy system of \(X\) and maps regularly onto the canonical residual
realization:
\[
\operatorname{Gen}_X(S)
\twoheadrightarrow
\mathsf D(S).
\]

\item If \(X\) is generated by \(S\) and behaviourally separated, then
\[
X\cong \mathsf D(S)
\]
as a Mealy system, and therefore
\[
\int_{\mathbb A}X
\cong
\operatorname{Syn}_{\mathbb A,\mathbb B}(S)
\]
as labelled input-discrete internal categories.

\item In the one-object Set case
\[
\mathbb A=M,
\qquad
\mathbb B=N,
\]
one has
\[
\mathsf D(S)
=
\{\partial_r\lambda\mid r\in M,\ \lambda\in S\}.
\]

\item In the one-object Set case, the syntactic transition-output monoid is
the image
\[
\operatorname{SynMon}_{M,N}(S)
=
\operatorname{im}
\left(
M
\to
\operatorname{MealyEnd}_N(\mathsf D(S))
\right),
\]
and is equivalently the quotient of \(M\) by the syntactic congruence
\[
{\equiv_S^{\operatorname{syn}}}.
\]

\item If
\[
M=I^\ast
\]
is free, then under the classical normal form
\[
\mathsf{Beh}_{I^\ast,N}\cong N^{I^+},
\]
the residual closure is
\[
\mathsf D(S)
=
\{\partial_w\beta\mid w\in I^\ast,\ \beta\in S\},
\]
where
\[
(\partial_w\beta)(u)=\beta(wu).
\]
The syntactic transition-output monoid is the image of
\[
I^\ast
\to
\operatorname{MealyEnd}_N(\mathsf D(S)).
\]
\end{enumerate}
\end{theorem}

\begin{proof}
Part (1) is Definition~\ref{def:image-specification}. Part (2) is
Definition~\ref{def:derivative-evaluation}. Part (3) is
Theorem~\ref{thm:residual-closure-operator}. Part (4) is
Corollary~\ref{cor:closed-behaviours-submealy}. Part (5) is
Theorem~\ref{thm:reflection-closed-behaviour-systems}. Part (6) is
Theorem~\ref{thm:canonical-nerode-quotient-presentation}. Part (7) is
Proposition~\ref{prop:residual-expression-is-mealy} and
Corollary~\ref{cor:residual-expression-quotient-mealy}. Part (8) is
Definition~\ref{def:syntactic-mealy-category} and
Definition~\ref{def:syntactic-mealy-category-labellings}. Part (9) is
Theorem~\ref{thm:syntactic-category-effective-quotient}. Part (10) is
Proposition~\ref{prop:generated-part-submealy} and
Proposition~\ref{prop:generated-part-to-residual-closure}. Part (11) is
Theorem~\ref{thm:generated-separated-realization-residual-closure}. Part
(12) is Proposition~\ref{prop:one-object-residual-closure}. Part (13) is
Definition~\ref{def:syntactic-transition-output-monoid} and
Theorem~\ref{thm:syntactic-monoid-quotient}. Part (14) is
Proposition~\ref{prop:free-input-residual-closure} and
Theorem~\ref{thm:free-input-syntactic-transition-output-monoid}.
\end{proof}